\section{Engines and Learning}\label{sec:engines}

CGT predicts which endgames should be hard for a chess engine. A number or a hot position shows up in the
material count or within the first few moves. An infinitesimal such as ${\uparrow}$, however, is a fight over one
spare tempo, and a nimber such as $\star$ is a fight over parity. Neither changes the material, so a static
evaluation sees nothing until the search reaches the end of the fight. In this section we test that prediction
on two engines, first on small positions that a tablebase can check and then on sums too large for any
tablebase. We end with a small learning experiment on Dawson's chess.

\subsection{A CGT Test Suite}

We built $233$ legal $8\times8$ positions with at most seven men, each labeled by the class of its CGT value:
a number, a switch $\game{x}{y}$ with $x>y$, a hot position containing a $\pm\BIG$ switch, an infinitesimal, a
nimber, or another value. Every position was checked against the Lichess tablebase~\cite{LichessTB}. In $181$
of them the winning moves predicted by CGT are confirmed by the tablebase, and in the other $52$ the model and
the tablebase disagree. We keep all $233$ and score the engines against the tablebase, since that is the real
answer.

\subsection{Stockfish and Lc0}

We ran Stockfish~19~\cite{Stockfish} on one thread with a $64$~MB hash and no tablebases. It solves all $233$
positions at $10^5$ nodes, and $232$ of them already at $10^3$ nodes \lab{C1}.

The interesting part is what happens at a small fixed depth (Table~\ref{tab:sf}). At depth~1 the engine reads
hot positions, numbers and switches correctly $73$--$81\%$ of the time. It reads infinitesimals correctly only
$30\%$ of the time at depth~1 and $26\%$ at depth~2, and nimbers $53\%$ of the time at depth~1 \lab{C1}. These
are exactly the tempo and parity fights that CGT singles out. Both classes reach $100\%$ by depth~8. This is one
engine on a small sample, so we treat it as a pilot result. Table~\ref{tab:sf} leaves out the seven positions
whose result is already decided.

\begin{table}[!t]
\centering
\caption{Stockfish~19 at Fixed Depth: Share of Evaluations (\%) That Match the Tablebase, by CGT Class}
\label{tab:sf}
\footnotesize
\begin{tabular}{@{}lrrrrrrr@{}}
\toprule
 & & \multicolumn{6}{c}{Search depth} \\
\cmidrule(l){3-8}
CGT class & $n$ & 1 & 2 & 4 & 6 & 8 & 12 \\
\midrule
hot ($\pm\BIG$ switch) & 68 & 81 & 78 & 91 & 97 & 99 & 100 \\
number & 57 & 79 & 81 & 84 & 93 & 95 & 98 \\
switch & 37 & 73 & 78 & 89 & 100 & 100 & 100 \\
other & 20 & 60 & 50 & 75 & 95 & 100 & 100 \\
nimber & 17 & 53 & 65 & 88 & 100 & 100 & 100 \\
infinitesimal & 27 & 30 & 26 & 59 & 85 & 100 & 100 \\
\bottomrule
\end{tabular}
\end{table}

Lc0~v0.32.1~\cite{Lc0}, with its bundled network, is perfect at $10^3$, $10^4$ and $10^5$ nodes: in every class,
every best move keeps the tablebase result and every evaluation matches it \lab{C1}. Since $216$ of the $233$
positions are wins, the $17$ draws are the real test, and Lc0 calls all of them correctly. So the suite is too
easy for both engines. A real contrast needs positions with more than seven men, where only CGT supplies the
answer.

\subsection{ChessSums}

ChessSums is a benchmark of pawn-endgame sums with exact answers. Each row is a sum of components drawn from a
library of $418$ components with $37$ distinct values. Every component has the same value under both
passed-pawn tests, all three race rules, and with en passant on or off, so the answer does not depend on our
conventions. The data set has $26{,}000$ rows. Every outcome was confirmed by a whole-board alpha--beta search,
and $4{,}887$ of them also by an independent minimax \lab{C2}. The test rows have more components than the
training rows, and some component types are held out entirely.

The benchmark also shows what the decomposition is worth. Alpha--beta on the whole board, which does not know
that the position splits, needs about $4$--$9$ times as many nodes for each component added \lab{C1}. Adding
up the CGT values of the components costs almost nothing.

\subsection{Engines on Sums}

To test the engines on ChessSums, we placed each row on an $8\times8$ board and added a frozen king group: a
tr\'ebuchet, or a shifted copy of it that our king-and-pawn solver confirms is a mutual zugzwang. The kings
then count as $0$ and the position is legal chess. These positions have $15$--$20$ men, far beyond any
tablebase. On $200$ of them, Stockfish~19 picks a move in the winning set only $27.5$--$28.0\%$ of the time,
from $10^3$ to $10^6$ nodes, and Lc0 does so $32.5$--$33.5\%$ of the time \lab{C1}. More nodes barely help
either engine, and both get worse as the number of components grows.

There is an important caveat. The ground truth here is the CGT model with frozen kings, not real chess. A move
that the model calls wrong may still win, or at least draw, at the board, because the convention can fail once
the kings take part (Section~\ref{sec:kings}). So these numbers measure how often the engines agree with the model; they do not prove
that the engines make real mistakes. Turning them into real-chess errors would need the kind of certification
that we carried out for $1/8$ in Section~\ref{sec:real}.

\subsection{Learning Dawson's Chess}

Can a neural network learn CGT? We tried the simplest case, normal-play Dawson's chess (Section~\ref{sec:dawson}).
There a position is a set of strips, and the first player wins exactly when the nim-sum of the Grundy values of the strips is not
zero~\cite{GuySmith1956,WinningWays}. We trained an LSTM value network, written in PyTorch~\cite{PyTorch}, on
positions of small total length and tested it on totals from $31$ to $80$. It reaches a test accuracy of
$0.82$, against about $0.79$ for always guessing the more common outcome \lab{C1}. Linear probes of the network
show why the margin is so small: its features tell whether a single strip is a first-player win, but they do
not recover the strip's Grundy value. The network has learned a hot/cold signal, not the nim-sum rule. We
report this as a negative result.
